% Notes-only D.93 candidate. Canonical checkpoint16 sources remain frozen.
\subsection{Unbounded finite compact kernels on actual ordinary fibres}

\begin{statement}{D.93}{Open}{Candidate: an ordinary-point family with explicit finite compact kernels}
For every integer \(r\ge1\), let
\[
 t_r=2^{5^r},\qquad
 s_r=\bigl(\kappa(t_r),\kappa(1-t_r)\bigr),
 \qquad n_r=5^r.                                      \tag{D.93.1}
\]
This is the abelian image of the actual rational-point section
of the tripod at \(t_r\). Retain the compact groups
\(\mathcal S_s\) and \(\mathcal K_s\) of D.63.
On the free \(\mathbb Z_5\)-module with basis
\(e_j\), \(j\in\mathbb Z/n_r\), define the two literal
permutation matrices
\[
 T_re_j=e_{j+1},\qquad
 \Pi_re_j=e_{2j},
 \qquad
 M_r=2(I+T_r)^3\Pi_r.                                 \tag{D.93.2}
\]
Then the entire separated compact kernel has the exact finite
presentation
\[
 \mathcal K_{s_r}\simeq\mathcal S_{s_r}
 \simeq
 \operatorname{coker}_{\mathbb Z_5}(M_r-I),
 \qquad
 \#\mathcal S_{s_r}=5^{\,1+r(r+1)/2}.                  \tag{D.93.3}
\]
In particular these finite orders are unbounded on actual
ordinary rational-point fibres. The first two integral Smith
calculations give
\[
 \mathcal K_{s_1}\simeq\mathbb Z/25,\qquad
 \mathcal K_{s_2}\simeq
        \mathbb Z/5\ \oplus\ \mathbb Z/125.             \tag{D.93.4}
\]
Equation (D.93.2) is a fully explicit finite integral matrix
for every \(r\). No general elementary-divisor pattern is
extrapolated from the first two cases.

Every indicated metabelian fibre is nonempty and has translation
group \(\prod_{\mathbb N}\mathbb Z_5\), with its compact product
topology. Exactly one ordinary rational point has its abelian
image, namely \(t_r\); no cusp has that abelian image. Thus every
other class in this metabelian fibre is neither a rational-point
class nor a cusp class. This is a statement about actual
geometrically pro-five metabelian sections. It does not classify
their higher-derived lifts or general abelian-pair existence.
\end{statement}

\begin{proof}
\paragraph{The global lattice and the exact local depths.}
Write \(B_r=2^{5^r}-1>0\). Since the sign has zero pro-five
Kummer class, the pair in (D.93.1) is
\[
        s_r=\bigl(5^r\kappa(2),\kappa(B_r)\bigr).
                                                               \tag{D.93.5}
\]
The integer \(B_r\) is odd and prime to five. It is not a
fifth power: put \(h=2^{5^{r-1}}\), so that
\((h-1)^5<h^5-1<h^5\). Consequently at least one prime
valuation of \(B_r\) is a unit in \(\mathbb Z_5\).
The first coordinate is supported only at two, and two does
not divide \(B_r\). Hence the actual closed away-five lattice is
\[
        L=5^r\mathbb Z_5(1,0)\oplus\mathbb Z_5(0,1),
             \qquad (a_5,b_5)=(0,0).                  \tag{D.93.6}
\]
This uses the complete valuation lattice, not only its reduction
modulo five. That reduction has rank one, so the earlier
residual-rank-two results do not give (D.93.3).

For a rational five-adic unit \(x\), use the integral coefficient
\[
                   \lambda_x=\frac{\log(x^4)}{4\log6}.
\]
It is the coefficient of \(\kappa(6)\) in the local unit class.
The denominator has valuation one. Direct reduction modulo
twenty-five gives
\[
 \lambda_2\equiv2\pmod5,\qquad
 2^{5^r}\equiv7\pmod{25},\qquad
 1-2^{5^r}\equiv-6\pmod{25},
 \qquad
 \lambda_{1-2^{5^r}}\equiv1\pmod5.                    \tag{D.93.7}
\]
Indeed \(5^r\equiv5\pmod{20}\), while two has order twenty
modulo twenty-five. Also
\(16-1\equiv15\pmod{25}\) and \(6^4-1\equiv20\pmod{25}\),
which give the two logarithmic congruences.
Put
\[
 \alpha=\lambda_2,\qquad
 \beta=\lambda_{1-2^{5^r}},\qquad
 \mathbf z=(5^r\alpha,\beta),\qquad
 \gamma=-\alpha/\beta\in\mathbb Z_5^\times .
                                                               \tag{D.93.8}
\]
These are the actual local coefficients in D.87.
The three global depths of \(a,b,a-b\) are respectively
\((r,0,0)\), and their three local depths are also \((r,0,0)\).
The first assertion follows from (D.93.6); the second follows
from (D.93.7), including the unit \(\beta-5^r\alpha\).
D.77 therefore gives
\[
                       \mathcal K_{s_r}\simeq
                       \mathcal S_{s_r}.              \tag{D.93.9}
\]
There is no discarded local free factor in this equality.

\paragraph{The actual image comparison and its natural map.}
Use \(\Lambda=\mathbb Z_5[[u,v]]\), \(t_1=1+u\),
\(t_2=1+v\), \(w=t_1t_2-1\), and \(P=uvw\).
Let \(I_0\) be the actual untwisted Alexander operator image,
including its cyclotomic coordinate \(\chi\), and let
\(N_0=\ker(\chi:I_0\to\mathbb Z_5^\times)\).
There is no valuation graph here, since (D.93.6) has zero
five-coordinate. D.87 supplies the actual common-complement
images
\[
              I_s=L\rtimes I_0,\qquad
              I_5=\mathbb Z_5\mathbf z\rtimes I_0 .
                                                               \tag{D.93.10}
\]
It also identifies the finite-localization kernel with its
natural coefficient map. In particular one finite coefficient
level suffices for all the local trivializations used in this
identification; no unrestricted locally trivial group is being
substituted for \(\mathcal S_s\).

Set \(q=(1+u)^{5^r}-1\). The local coefficient quotient and
translation ideal are
\[
 R=\Lambda/(t_1^{5^r\alpha}t_2^\beta-1),
                    \qquad J=(q,v)R.                 \tag{D.93.11}
\]
Since \(\beta\) is a unit, the first relation eliminates \(t_2\)
integrally. Thus, as actual topological rings,
\[
 R=\mathbb Z_5[[u]],\qquad
 t_2=t_1^{5^r\gamma},\qquad
 v=(1+q)^\gamma-1,\qquad J=qR.                         \tag{D.93.12}
\]
The series \(v/q\) has unit constant coefficient \(\gamma\).
Likewise
\[
 w=t_1^{1+5^r\gamma}-1
\]
is \(u\) times a unit, since \(1+5^r\gamma\) is a unit.
Consequently the ideal generated by the actual \(P\) in \(R\) is
\[
                              PR=u^2qR.               \tag{D.93.13}
\]
No rationalization or splitting into cyclotomic characters
has been used in these integral identities.

For every \(g\in N_0\), the exact beta boundary identities
of D.39 give \(\mathcal B_g-1\in(P)\). The primitive actual
local cubic element in D.77 has
\(\mathcal B_\delta-1=P E_\delta\), with \(E_\delta\) a unit.
Hence the ideal of all these actual operator differences is
exactly \((P)\), on \(R\) and on the invariant ideal \(J\).
Ordinary duality applied to D.87.4 therefore gives
\[
 \mathcal S_s=
 \ker\left\{
       (J/PJ)(-1)_\Gamma\longrightarrow
       (R/PR)(-1)_\Gamma\right\},
                    \qquad \Gamma=\mathbb Z_5^\times .
                                                               \tag{D.93.14}
\]
The map is still induced by the actual inclusion \(J\subset R\).
The twist \((-1)\) records precisely the Cartier-to-ordinary
duality convention. These are compact coinvariants, with closed
images, and (D.93.14) is a map of compact modules.

\paragraph{Retain the generator before taking the two low jets.}
Scalar division by the retained generator \(q\) identifies
\(J/PJ\) with
\[
                         T=R/(u^2q).
\]
On this scalar module the actual Cartier-untwisted action is
\[
        \chi(g)^{-1}S_{\chi(g)}(t_1^{5^r})\rho_0(g),
          \qquad S_c(t)=\frac{t^c-1}{t-1},             \tag{D.93.15}
\]
where \(\rho_0\) includes its substitution \(t_i\mapsto
t_i^{\chi(g)}\). The natural map in (D.93.14) is multiplication
by \(q\), not the identity on two isomorphic scalar quotients.

The quotient \(T/u^2T=R/(u^2)\) is free of rank two over
\(\mathbb Z_5\). Its successive actual characters are
\(\chi^2,\chi^3\): the Alexander multiplier has augmentation
\(\chi^2\), the retained \(q\) contributes \(\chi\), the
Cartier untwist removes one factor \(\chi\), and the linear
substitution supplies the next factor. Both characters are
nontrivial on the order-four tame subgroup. Exact averaging
by that subgroup shows that every continuous
\(\Gamma\)-homology group of this two-jet quotient is zero.
The compact exact homology sequence consequently gives an
isomorphism
\[
                 (u^2T)_\Gamma\xrightarrow{\ \sim\ }T_\Gamma.
                                                               \tag{D.93.16}
\]
This proves both surjectivity and injectivity; no claim that a
coinvariant functor is generally left exact is needed.

Since \(R\) is a domain, scalar division by \(u^2\) gives
\[
                u^2T\simeq R/(q)
                    \simeq\mathbb Z_5[C_{5^r}].        \tag{D.93.17}
\]
On this quotient \(t_1^{5^r}=t_2=1\).
The exact untwisted Alexander action is
\[
 \rho_0(g)F=
 S_\chi(t_1)S_\chi(t_2)\mathcal B_g(t_1,t_2)
                 F(t_1^\chi,t_2^\chi).
\]
The boundary identity \(\mathcal B_g(t,1)=1\) is exact for
every multiplier. Combining it with (D.93.15) and the
substitution on the retained \(u^2\) gives on (D.93.17)
\[
                       F(t)\longmapsto
                  \chi\,S_\chi(t)^3 F(t^\chi).         \tag{D.93.18}
\]
Thus neither an unspecified Frobenius constant nor an
approximate weight action remains.

Multiplication by \(q\) sends the actual \(u^2\)-layer to
zero in \(R/PR\), since \(u^2q=0\) there. By (D.93.16)
the entire natural map in (D.93.14) is zero.
It follows that
\[
             \mathcal S_{s_r}\simeq
       \bigl(\mathbb Z_5[C_{5^r}],
                  \chi S_\chi(t)^3\sigma_\chi\bigr)_\Gamma .
                                                               \tag{D.93.19}
\]
The element two topologically generates \(\mathbb Z_5^\times\):
its order modulo five is four and
\(v_5(2^4-1)=1\). Its actual action in (D.93.18) is precisely
the matrix \(M_r\) of (D.93.2). The image of \(M_r-I\)
on this finite free compact module is closed. Density of the
cyclic subgroup generated by two therefore identifies its
cokernel with the full compact \(\Gamma\)-coinvariants.
This proves the finite presentation in (D.93.3).

\paragraph{The all-depth determinant.}
Compute the determinant after adjoining the \(5^r\)-th roots
of unity in characteristic zero; this does not change it.
The evaluation at \(t=1\) has scalar \(16\).
For each \(1\le j\le r\), multiplication of the exponents by
two is one orbit on the primitive \(5^j\)-th roots, of length
\(\ell_j=4\cdot5^{j-1}\). On that orbit the product of the
weights in (D.93.18) is
\[
 2^{\ell_j}\prod_{k=0}^{\ell_j-1}
       (1+\zeta^{2^k})^3
 =2^{\ell_j},
 \qquad
 1+\zeta^{2^k}
     =\frac{\zeta^{2^{k+1}}-1}{\zeta^{2^k}-1}.          \tag{D.93.20}
\]
The product telescopes because the orbit closes and none
of these primitive roots is one. For a weighted permutation
cycle the determinant of identity minus the cycle is one
minus its weight product. Consequently
\[
 \det(I-M_r)
      =(1-16)\prod_{j=1}^r
                 (1-2^{\,4\cdot5^{j-1}}).             \tag{D.93.21}
\]
All factors are nonzero. The elementary lifting-the-exponent
calculation gives
\[
 v_5(16-1)=1,\qquad
 v_5(2^{\,4\cdot5^{j-1}}-1)=j.
 \]
For example this follows inductively by expanding
\((1+5c)^5-1\), which increases the valuation by exactly
one when \(c\) is a unit. Thus the cokernel is finite of the
order in (D.93.3), with its finite compact topology.
The exact literal integral Smith forms for \(r=1,2\) are
executed in the accompanying certificate and give (D.93.4).
Those two executions are not used to prove (D.93.21).

\paragraph{The actual section fibres and their ordinary points.}
The rational point \(t_r\) supplies a full section, hence an
actual metabelian class above \(s_r\). The pair has finite
prime support. D.74 therefore identifies the translation
group of its metabelian fibre with
\(\prod_{\mathbb N}\mathbb Z_5\).
The two Kummer coordinates are linearly independent over
\(\mathbb Q_5\), by their disjoint nonempty prime supports.
They lie on none of the three completed cusp lines in D.35.

For a nonzero rational number \(x\), its pro-five Kummer
class vanishes exactly when \(x=\pm1\); this follows directly
from its prime valuations. If an ordinary rational point
\(z\) has abelian pair \(s_r\), it follows that
\[
                  z=\pm t_r,\qquad
                  1-z=\pm(1-t_r).
 \]
The choice \(z=-t_r\) makes the second equality impossible
in either sign. Hence \(z=t_r\), as claimed. These statements
classify the ordinary point images in this fibre, not the
full section classes of the geometric fundamental group.
\end{proof}
